% GENERATED FILE. Edit canonical lessons, then run npm run generate:books.
% canonical-source-hash: fnv1a64:2394ebd9ad3389b7
% canonical-lessons: AR-C505-jumjuma, AR-C505-dil, AR-C505-fakhidh, AR-C505-wirk, AR-C505-mafsil

\chapter{Bones and joints}
\label{ch:ar-bones-and-joints}
\hlchaptermodality{505}

\begin{chapteropening}
\textbf{By the end of this chapter:} \emph{I can name the skull, a rib, a thigh, a hip and a joint.}

Complete the last lesson of chapter 505: I can name the skull, a rib, a thigh, a hip and a joint.
\end{chapteropening}

\section[\texorpdfstring{\ar{جمجمة} (jumjuma)}{جمجمة (jumjuma)}]{\ar{جمجمة} (jumjuma) — a skull}

\label{lesson:AR-C505-jumjuma}



\begin{quote}
Before the new one: say the Arabic for the intestines, then the Arabic for the pulse.
\end{quote}

\subsection*{You'll want to know: \ar{جمجمة}}
\textbf{\ar{جمجمة}} — \emph{jumjuma} — \textquotedblleft{}a skull\textquotedblright{}.

\begin{quote}
\ar{الجمجمة} \ar{تحمي} \ar{الدماغ}.
\end{quote}

\begin{quote}
\emph{al-jumjuma taḥmī ad-dimāgh.}
\end{quote}

\begin{quote}
The skull protects the brain.
\end{quote}

Plural: \ar{جماجم}, \emph{jamājim}.

\begin{cousinweb}
\textbf{\ar{جمجمة}} (\emph{jumjuma}) is a native Arabic word on a doubled four-letter shape, \emph{j-m-j-m}. EtymDB sets it beside Hebrew \emph{gulgolet} and Aramaic \emph{gulgaltā}, \textquotedblleft{}a skull\textquotedblright{}, but how they are related is not settled. Its plural \emph{jamājim} follows the usual four-letter pattern.
\end{cousinweb}

\begin{grammarlens}[title={What you've built}]
One more word for describing what you lived and felt.
\end{grammarlens}

\subsection*{Your turn}
\begin{itemize}
\raggedright
  \item \emph{Say it:} \emph{jumjuma}
  \item \emph{Say it:} \emph{jumjuma}, once more
  \item \emph{Say it:} \emph{al-jumjuma taḥmī ad-dimāgh.} — The skull protects the brain.
  \item \emph{Recall:} say the Arabic for the intestines, then the Arabic for the pulse, then say \emph{jumjuma} again
\end{itemize}

\begin{tcolorbox}[breakable,colback=teal!4,colframe=teal!35!black,title={Before you move on}]
What does \ar{جمجمة} mean? (\textquotedblleft{}A skull\textquotedblright{}.) Say it once more.
\end{tcolorbox}
\section[\texorpdfstring{\ar{ضلع} (ḍilʿ)}{ضلع (ḍilʿ)}]{\ar{ضلع} (ḍilʿ) — a rib}

\label{lesson:AR-C505-dil}



\begin{quote}
Before the new one: say the Arabic for the pulse, then the Arabic for a skull.
\end{quote}

\subsection*{You'll want to know: \ar{ضلع}}
\textbf{\ar{ضلع}} — \emph{ḍilʿ} — \textquotedblleft{}a rib\textquotedblright{}.

\begin{quote}
\ar{كسر} \ar{اللاعب} \ar{ضلعا} \ar{في} \ar{المباراة}.
\end{quote}

\begin{quote}
\emph{kasara al-lāʿib ḍilʿan fī al-mubārā.}
\end{quote}

\begin{quote}
The player broke a rib in the match.
\end{quote}

Plural: \ar{أضلاع}, \emph{aḍlāʿ}.

In geometry \ar{ضلع} is also the side of a triangle.

\begin{cousinweb}
\textbf{\ar{ضلع}} (\emph{ḍilʿ}) is built on the native Arabic root \emph{ḍ-l-ʿ}, \textquotedblleft{}a rib; to be crooked\textquotedblright{} (Lane). Lane links the rib to the sense of being bent.
\end{cousinweb}

\begin{grammarlens}[title={What you've built}]
One more word for describing what you lived and felt.
\end{grammarlens}

\subsection*{Your turn}
\begin{itemize}
\raggedright
  \item \emph{Say it:} \emph{ḍilʿ}
  \item \emph{Say it:} \emph{ḍilʿ}, once more
  \item \emph{Say it:} \emph{kasara al-lāʿib ḍilʿan fī al-mubārā.} — The player broke a rib in the match.
  \item \emph{Recall:} say the Arabic for the pulse, then the Arabic for a skull, then say \emph{ḍilʿ} again
\end{itemize}

\begin{tcolorbox}[breakable,colback=teal!4,colframe=teal!35!black,title={Before you move on}]
What does \ar{ضلع} mean? (\textquotedblleft{}A rib\textquotedblright{}.) Say it once more.
\end{tcolorbox}
\section[\texorpdfstring{\ar{فخذ} (fakhidh)}{فخذ (fakhidh)}]{\ar{فخذ} (fakhidh) — a thigh}

\label{lesson:AR-C505-fakhidh}



\begin{quote}
Before the new one: say the Arabic for a skull, then the Arabic for a rib.
\end{quote}

\subsection*{You'll want to know: \ar{فخذ}}
\textbf{\ar{فخذ}} — \emph{fakhidh} — \textquotedblleft{}a thigh\textquotedblright{}.

\begin{quote}
\ar{أكلنا} \ar{فخذ} \ar{الدجاج} \ar{المشوي}.
\end{quote}

\begin{quote}
\emph{akalnā fakhidh ad-dajāj al-mashwī.}
\end{quote}

\begin{quote}
We ate the grilled chicken thigh.
\end{quote}

Plural: \ar{أفخاذ}, \emph{afkhādh}.

\begin{cousinweb}
\textbf{\ar{فخذ}} (\emph{fakhidh}) is built on the native Arabic root \emph{f-kh-dh}, \textquotedblleft{}the thigh\textquotedblright{} (Lane). Lane also gives fakhidh as a branch of a tribe, below the clan.
\end{cousinweb}

\begin{grammarlens}[title={What you've built}]
One more word for describing what you lived and felt.
\end{grammarlens}

\subsection*{Your turn}
\begin{itemize}
\raggedright
  \item \emph{Say it:} \emph{fakhidh}
  \item \emph{Say it:} \emph{fakhidh}, once more
  \item \emph{Say it:} \emph{akalnā fakhidh ad-dajāj al-mashwī.} — We ate the grilled chicken thigh.
  \item \emph{Recall:} say the Arabic for a skull, then the Arabic for a rib, then say \emph{fakhidh} again
\end{itemize}

\begin{tcolorbox}[breakable,colback=teal!4,colframe=teal!35!black,title={Before you move on}]
What does \ar{فخذ} mean? (\textquotedblleft{}A thigh\textquotedblright{}.) Say it once more.
\end{tcolorbox}
\section[\texorpdfstring{\ar{ورك} (wirk)}{ورك (wirk)}]{\ar{ورك} (wirk) — a hip}

\label{lesson:AR-C505-wirk}



\begin{quote}
Before the new one: say the Arabic for a rib, then the Arabic for a thigh.
\end{quote}

\subsection*{You'll want to know: \ar{ورك}}
\textbf{\ar{ورك}} — \emph{wirk} — \textquotedblleft{}a hip\textquotedblright{}.

\begin{quote}
\ar{تشعر} \ar{جدتي} \ar{بألم} \ar{في} \ar{الورك}.
\end{quote}

\begin{quote}
\emph{tashʿuru jaddatī bi-alam fī al-wirk.}
\end{quote}

\begin{quote}
My grandmother feels a pain in her hip.
\end{quote}

\begin{cousinweb}
\textbf{\ar{ورك}} (\emph{wirk}) is built on the native Arabic root \emph{w-r-k}, \textquotedblleft{}the hip, the haunch\textquotedblright{} (Lane). The word is old and common to the language.
\end{cousinweb}

\begin{grammarlens}[title={What you've built}]
One more word for describing what you lived and felt.
\end{grammarlens}

\subsection*{Your turn}
\begin{itemize}
\raggedright
  \item \emph{Say it:} \emph{wirk}
  \item \emph{Say it:} \emph{wirk}, once more
  \item \emph{Say it:} \emph{tashʿuru jaddatī bi-alam fī al-wirk.} — My grandmother feels a pain in her hip.
  \item \emph{Recall:} say the Arabic for a rib, then the Arabic for a thigh, then say \emph{wirk} again
\end{itemize}

\begin{tcolorbox}[breakable,colback=teal!4,colframe=teal!35!black,title={Before you move on}]
What does \ar{ورك} mean? (\textquotedblleft{}A hip\textquotedblright{}.) Say it once more.
\end{tcolorbox}
\section[\texorpdfstring{\ar{مفصل} (mafṣil)}{مفصل (mafṣil)}]{\ar{مفصل} (mafṣil) — a joint}

\label{lesson:AR-C505-mafsil}



\begin{quote}
Before the new one: say the Arabic for a thigh, then the Arabic for a hip.
\end{quote}

\subsection*{You'll want to know: \ar{مفصل}}
\textbf{\ar{مفصل}} — \emph{mafṣil} — \textquotedblleft{}a joint\textquotedblright{}.

\begin{quote}
\ar{عند} \ar{جدي} \ar{ألم} \ar{في} \ar{المفاصل}.
\end{quote}

\begin{quote}
\emph{ʿinda jaddī alam fī al-mafāṣil.}
\end{quote}

\begin{quote}
My grandfather has pain in his joints.
\end{quote}

Plural: \ar{مفاصل}, \emph{mafāṣil}.

\begin{cousinweb}
\textbf{\ar{مفصل}} (\emph{mafṣil}) is built on the native Arabic root \emph{f-ṣ-l}, \textquotedblleft{}to separate, to divide\textquotedblright{} (Lane). Faṣl, a season, which you know, shares the root: both name a division. A mafṣil is where the body divides.
\end{cousinweb}

\begin{grammarlens}[title={What you've built}]
One more word for describing what you lived and felt.
\end{grammarlens}

\subsection*{Your turn}
\begin{itemize}
\raggedright
  \item \emph{Say it:} \emph{mafṣil}
  \item \emph{Say it:} \emph{mafṣil}, once more
  \item \emph{Say it:} \emph{ʿinda jaddī alam fī al-mafāṣil.} — My grandfather has pain in his joints.
  \item \emph{Recall:} say the Arabic for a thigh, then the Arabic for a hip, then say \emph{mafṣil} again
\end{itemize}

\begin{tcolorbox}[breakable,colback=teal!4,colframe=teal!35!black,title={Before you move on}]
What does \ar{مفصل} mean? (\textquotedblleft{}A joint\textquotedblright{}.) Say it once more.
\end{tcolorbox}
